\section*{الفصل \ref{ch:b2:mammal-reproduction} --- التكاثر الجنسي عند الثدييات}

\begin{solution}{exo:b2:mammal-reproduction:1}
تنقسم المنسليات من البلوغ مدى الحياة؛ أما الخلايا المبيضية فلا
تنقسم إلا في الجنين. ويعطي \omterm{def:b2:meiosis-heredity:meiosis}{التقسم المنصّف} الواحد أربع نطاف لكن
بويضة واحدة (وثلاثة أجسام قطبية). وتستغرق النطفة نحو 64 يومًا زائد
أسبوعين من النضج؛ أما البويضة فموقوفة في الطور التمهيدي الأول من
قبل الولادة إلى شهر إباضتها --- أي عقودًا --- وموقوفة ثانيةً في
الطور الاستوائي الثاني حتى الإلقاح. ويصنع الرجل $10^{8}$ نطفة في
اليوم، أي نحو $10^{12}$ في العمر؛ وتبيض المرأة نحو 400 بويضة.
\end{solution}

\begin{solution}{exo:b2:mammal-reproduction:2}
الارتباط \omterm{prop:b2:mammal-reproduction:fertilisation}{بالمنطقة الشفافة}؛ \omterm{prop:b2:mammal-reproduction:fertilisation}{وتفاعل الجسيم القمي} وهضم
ممر؛ واندماج الغشاءين ودخول نواة النطفة؛
وموجة الكالسيوم \omterm{prop:b2:mammal-reproduction:fertilisation}{والتفاعل القشري} (تصلّب المنطقة، وفقد
مستقبلاتها: أي الحاجز أمام \omterm{thm:b2:mammal-reproduction:poisson}{تعدد النطاف})؛ وإتمام \omterm{def:b2:meiosis-heredity:meiosis}{التقسم المنصّف}
الثاني وطرح الجسم القطبي الثاني؛ وتكوّن النواتين
الأوليتين، ونسخ الدنا، وأول تقسم خيطي.
\end{solution}

\begin{solution}{exo:b2:mammal-reproduction:3}
\omterm{def:b2:mammal-reproduction:implantation}{الأرومة المغذية} (الطبقة الخارجية) تصنع \omterm{def:b2:mammal-reproduction:placenta}{المشيمة} والأغشية؛ \omterm{def:b2:mammal-reproduction:implantation}{والكتلة
الخلوية الداخلية} تصنع الجنين. ويبدأ \omterm{def:b2:mammal-reproduction:implantation}{الانغراس} في اليوم 6 إلى 7
بعد الإلقاح، في بطانة رحمية هيأها البروجستيرون
من \omterm{def:b2:mammal-reproduction:ovary}{الجسم الأصفر}.
\end{solution}

\begin{solution}{exo:b2:mammal-reproduction:4}
الأكسجين: بالانتشار. والغلوكوز: بنقل ميسَّر. والكريات الحمر
الأمومية: لا تعبر (فالدمان لا يختلطان أبدًا). وIgG: بنقل
بوساطة مستقبل. والبولة: بالانتشار، من الجنين إلى الأم. والكحول:
ينتشر بحرية. والبكتيريا: لا تعبر في معظمها (وبعضها، كالإفرنجي
والليستيريا، يعبر).
\end{solution}

\begin{solution}{exo:b2:mammal-reproduction:5}
$m = 3, 1, 0.4, 0.1$: ومنه $P = 0.95, 0.63, 0.33, 0.095$. ويهبط دون
النصف حين $m < \ln 2 = 0.69$، أي دون $7\times 10^{7}$ نطفة. والمنحني
يتشبع: ففوق نحو $10^{8}$ لا تضيف النطاف الزائدة شيئًا يُذكر،
ودونه تهبط الخصوبة بنسبة العدد تقريبًا.
\end{solution}

\begin{solution}{exo:b2:mammal-reproduction:6}
قبل البلوغ: $7\times 10^{5}$ بويضة تُفقد في 13 سنة (\num{4750} يومًا):
أي نحو 150 في اليوم. وبعده: \num{299000} تُفقد في 37 سنة (\num{13500}
يومًا): أي 22 في اليوم؛ \omterm{def:b2:mammal-reproduction:ovary}{والإباضة} $400/\num{13500} = 0.03$ في اليوم.
فالرتق يفوق \omterm{def:b2:mammal-reproduction:ovary}{الإباضة} سبعمئة إلى واحد؛ ولا يُنفق المخزون
\omterm{def:b2:mammal-reproduction:ovary}{بالإباضة} بل بالموت.
\end{solution}

\begin{solution}{exo:b2:mammal-reproduction:7}
البالغ: $(30/27)^{2.7} = 1.33$، ومنه $S = 1.33/2.33 = 0.57$. والجنين:
$(30/19)^{2.7} = 3.4$، ومنه $S = 3.4/4.4 = 0.77$. فعند الضغوط
المنخفضة في \omterm{def:b2:mammal-reproduction:placenta}{المشيمة} يحمّل الدم الجنيني خُمسًا من الأكسجين أكثر مما
كان ليحمّله دم بالغ، ويفرّغه في الأنسجة عند ضغوط يكون فيها
هيموغلوبين البالغ نصف فارغ سلفًا.
\end{solution}

\begin{solution}{exo:b2:mammal-reproduction:8}
$2^{k} = 25$: ومنه $k = 4.6$ مضاعفة، أي \qty{9.3}{d}: فاليوم 16 إلى 17 بعد
الإلقاح؛ واليوم 30 إلى 31 بعد آخر دورة --- أي نحو اليوم
الذي تتأخر فيه الدورة.
\end{solution}

\begin{solution}{exo:b2:mammal-reproduction:9}
$3\times 10^{-8}\times 6\times 10^{4}\times 20/10^{-3} =
\qty{36}{mL/min}$ في مقابل طلب قدره \qty{25}{mL/min}: فيهبط الهامش
من ثمانية إلى واحد ونصف؛ ويصير الجنين ناقص الأكسجة تحت
أي حمل إضافي، وينمو ببطء وقد يحتاج إلى ولادة مبكرة.
\end{solution}

\begin{solution}{exo:b2:mammal-reproduction:10}
(أ) تتكون الخصيتان وتفرزان الهرمونين معًا؛ فيزيل \omterm{prop:b2:mammal-reproduction:sex}{الهرمون المضاد
لمولر} القنوات المولرية، لكن التستوستيرون لا يستطيع العمل: فتضمر
القنوات الولفية وتكون الأعضاء التناسلية الخارجية أنثوية --- أي
امرأة بخصيتين وبلا رحم. (ب) يعمل التستوستيرون ولا يعمل \omterm{prop:b2:mammal-reproduction:sex}{الهرمون
المضاد لمولر}: فذكر ببربخ وأسهر وأعضاء تناسلية ذكرية، مع رحم
وقناتي بيض. (ج) مبيضان ومشتقات مولرية (فلا خصية ولا هرمون مضاد
لمولر)؛ وتُذكّر أندروجينات الكظر الأعضاء التناسلية الخارجية بدرجات
متفاوتة. (د) تصنع \emph{SRY} خصيتين تفرضان المسلك الذكري: أي رجل
بصبغيي X، عقيم لأن مورّثات Y الخاصة \omterm{def:b2:mammal-reproduction:testis}{بتكوّن النطاف}
غائبة.
\end{solution}

\begin{solution}{exo:b2:mammal-reproduction:11}
الأشيع (\qty{68}{\%}) هو الانشطار في اليوم 4 إلى 8: أي بعد
تكوّن \omterm{def:b2:mammal-reproduction:implantation}{الأرومة المغذية} (\omterm{def:b2:mammal-reproduction:placenta}{فمشيمة} واحدة) وقبل السلى
(فكيسان). أما قبل اليوم 3 فيصنع كل نصف أرومته المغذية،
ومن ثم مشيمتان (\qty{30}{\%})؛ وبعد اليوم 8 يكون السلى
قد وُضع سلفًا فيكون مشتركًا (\qty{2}{\%}). فانشطار \omterm{def:b2:mammal-reproduction:implantation}{الكتلة
الخلوية الداخلية} بعد التزام \omterm{def:b2:mammal-reproduction:implantation}{الأرومة المغذية} هو بيّنًا
أسهل الأحداث.
\end{solution}

\begin{solution}{exo:b2:mammal-reproduction:12}
يستثمر الجرابي قليلًا قبل الولادة: حمل قصير، وصغير
ضئيل، وإرضاع يمكن قطعه بإسقاط صغير الجراب
إن جاء الجفاف، فيكون خطر الأم في المحاولة الواحدة قليلًا
وتستطيع أن تعيد المحاولة حالًا. أما المشيمي فيستثمر بكثافة في
حمل طويل بصغير كبير حسن النمو ينجو أفضل
عند الولادة، بثمن أم ملتزمة أشهرًا
ومعرّضة في تلك الأثناء. وفي البيئات المنتجة المتوقعة يغلب
نجاءُ صغار المشيمي الأعلى؛ أما مناخ أستراليا المتقلب
فيكافئ خيار الجرابي في التخلي، وعزلتها
أبقت المشيمية خارجًا حتى عهد قريب.
\end{solution}

\begin{solution}{pb:b2:mammal-reproduction:1}
\textbf{1.} $m = 3\times 10^{8}\times 10^{-8} = 3$؛ ومنه $P = 1 - e^{-3} =
0.95$.
\textbf{2.} $1 - e^{-3}(1 + 3) = 0.80$: فلولا الحاجز لكانت أربع
إلقاحات من كل خمس متعددة النطاف والجنين ثلاثي الصيغة.
\textbf{3.} $0.15/5\times 10^{-5} = \qty{3000}{s}$، أي خمسون دقيقة؛
وتحملها تقلصات الرحم وقناة البيض.
\textbf{4.} $300/3\times 10^{8} = 10^{-6}$: وتضيع البقية في المهبل
الحمضي، وفي مخاط عنق الرحم، وأمام البالعات في الرحم، وفي قناة
البيض الخطأ.
\textbf{5.} من ثلاثة أيام قبل \omterm{def:b2:mammal-reproduction:ovary}{الإباضة} إلى يوم بعدها: أي نحو
أربعة أيام.
\textbf{6.} $m = 0.5$، و$P = 0.39$؛ والدورات المنتظرة $1/P = 2.5$.
\textbf{7.} متقدرات النطفة القليلة، في القطعة الوسطى، تُوسَم
وتُتلف بعد الدخول؛ أما البويضة فتحمل مئة ألف.
\textbf{8.} $120/16 = 7.5$: أي سبعة انقسامات، و$2^{7} = 128$ خلية ---
أي نحو مئة.
\textbf{9.} $\log_2 25 = 4.6$ مضاعفة، أي \qty{9.3}{d}: فاليوم 16 إلى 17؛
واليوم 30 إلى 31 بعد الدورة.
\textbf{10.} يضمر \omterm{def:b2:mammal-reproduction:ovary}{الجسم الأصفر} بعد نحو اثني عشر يومًا من
\omterm{def:b2:mammal-reproduction:ovary}{الإباضة} ما لم ينقذه hCG؛ ودونه يهبط البروجستيرون، وتنسلخ
البطانة ويضيع الجنين مع الدورة.
\textbf{11.} $\log_2 10^{5} = 16.6$ مضاعفة؛ وعند الأسبوع العاشر تصنع
\omterm{def:b2:mammal-reproduction:placenta}{المشيمة} بروجستيرونها بنفسها ولا يعود \omterm{def:b2:mammal-reproduction:ovary}{الجسم الأصفر}
لازمًا.
\textbf{12.} تجلس البويضات في الطور التمهيدي الأول عقودًا وتتحلل
الكوهيزينات الماسكة لثنائيات التكافؤ، فيرتفع عدم الانفصال في
\omterm{def:b2:meiosis-heredity:meiosis}{التقسم المنصّف} الأول ارتفاعًا حادًا مع عمر الأم؛ أما النطاف فتُصنع
طازجة في شهرين.
\textbf{13.} بعد \omterm{def:b2:mammal-reproduction:implantation}{الأرومة المغذية} (اليوم 5) وقبل السلى (اليوم
8): أي \omterm{def:b2:mammal-reproduction:placenta}{مشيمة} واحدة وكيسان.
\textbf{14.} $3\times 10^{-8}\times 1.2\times 10^{5}\times 20/3.5\times
10^{-4} = \qty{206}{mL/min}$.
\textbf{15.} $7\times 3.5 = \qty{24.5}{mL/min}$: أي هامش قدره ثمانية.
\textbf{16.} $500\times 0.16 = \qty{80}{mL/min}$؛ والاستخلاص
$24.5/80 = \qty{31}{\%}$.
\textbf{17.} $5\times 3.5\times 1440 = \qty{25}{g/d}$؛ وهو في \qty{28}{L}
من الدم الأمومي؛ والجريان اليومي \qty{720}{L}: فيأخذ الجنين نحو
\qty{4}{\%} من الغلوكوز العابر.
\textbf{18.} هيموغلوبين الجنين مشبع بنسبة \qty{77}{\%} عند
\qty{30}{mmHg}، وللجنين كريات حمر أكثر ونتاج قلبي
عالٍ، وأنسجته متكيفة للعمل عند ضغط منخفض --- فهو
يعيش، في الواقع، على قمة إيفرست.
\textbf{19.} التبادل الغازي بلا رئة عاملة، والتغذية
والإطراح بلا معي عامل ولا كلية، والهرمونات التي
تصون الحمل؛ ولا تستطيع صدّ كل ذيفان --- فالكحول
والنيكوتين وبعض الفيروسات تعبر.
\textbf{20.} تمدد عنق الرحم $\to$ الوطاء $\to$ الأوكسيتوسين $\to$
التقلص $\to$ تمدد أكثر، حتى تزيل الولادة المنبه.
وقبل التمام يكون للرحم الذي يسوده البروجستيرون مستقبلات
أوكسيتوسين قليلة ولا وصلات فجوية، فلا تنتشر التقلصات ولا
تستطيع الحلقة أن تنغلق.
\textbf{21.} $750\times 2.9 = \qty{2.2}{MJ}$؛ والكلفة $2.2/0.8 =
\qty{2.7}{MJ}$.
\textbf{22.} النمو $30\times 6 = \qty{0.18}{MJ}$، والصيانة
$300\times 4 = \qty{1.2}{MJ}$: أي \qty{1.4}{MJ}، نحو ثلثي
طاقة الحليب؛ ويذهب الباقي إلى النشاط والدفء والفواقد.
\textbf{23.} $270\times 2.7 = \qty{730}{MJ}$، أي ضعفي كلفة الحمل
ونصف.
\textbf{24.} ترفع الرضاعة البرولاكتين وتكبت الإطلاق النبضي
للهرمون الذي يقود \omterm{def:b2:mammal-reproduction:ovary}{الإباضة}، فلا تعود \omterm{def:b2:mammal-reproduction:ovary}{الإباضة}
إلا بعد أشهر؛ ودون وسائل منع الحمل يباعد هذا بين الولادات سنتين
إلى ثلاث.
\textbf{25.} $P = 0.95$؛ والاختبار يصير موجبًا في اليوم 17 بعد الإلقاح
(أي اليوم 31 من الدورة)؛ وهامش الأكسجين ثمانية؛ وحليب يوم يكلف
الأم \qty{2.7}{MJ}.
\end{solution}
